% sec-01: the signal.  Table 1, Figure 1 (mermaid-type flowchart).
\section{The Signal, and Why It Is Spent First}

\subsection{What arrived}

Two LinkedIn messages, a few days apart, from two people who do not know each
other and did not know the other had written. Each asked, in their own words,
whether ChemicalQDevice was hiring. Neither was answering a posting, because
there was none: the company has published no job listing, no careers page and
no request for applicants, and it has used no recruiter. Each message named the
company's public work as the reason for writing.

\subsection{The standing fact behind every day of this block}

On August 26, 2026 the FDA approved Rasonque (daraxonrasib), Revolution
Medicines' RAS(ON) multi-selective inhibitor, as a first-in-class targeted
therapy for metastatic pancreatic adenocarcinoma after prior systemic therapy
\cite{fdarasonque2026,fdadisco2026,revmedapproval2026}. The approval rests on the
Phase 3 RASolute 302 trial, which reported a median overall survival of 13.2
months against 6.6 months for chemotherapy in the RAS G12 population
\cite{rasolute302,ctgov302}, and it followed the developer's Phase 1/2 program,
registered as the Phase 1/1b study of RMC-6236 \cite{ctgov001}.

ChemicalQDevice's AI papers were developed independently of that program and
simultaneously with it. In June 2025, while RASolute 302 was enrolling, the
company ranked a daraxonrasib combination first of forty pancreatic cancer
meta-analyses \cite{daraxident}. In August 2025 its ten-arm quantitative systems
pharmacology simulation returned 12.8 months against 5.4 for chemotherapy
\cite{qspmetpancre}. Between June and August 2026, before and after the approval,
it deposited a Phase 1 protocol, an investigational new drug application, a
Phase 2 protocol and two funding applications around the molecule
\cite{onpremwhippl,phase1ind,phase2proto,fundapp1,fundapp2}. Each carries a date,
and the chronology is kept in one public file \cite{daraxstory2026}. That record
is what the two applicants read.

\subsection{Why two inquiries are evidence}

Federal small-business review scores the team. Until last week the company's
answer was a plan: eleven roles at 3.95 award-funded full-time equivalents and a
\$521,000 personnel line inside a \$700,000 annual direct cost \cite{movein2026}.
Table 1 sets out what two unprompted inquiries add to that plan, and, in the
same row, what they cannot show.

\begin{apptable}
\begin{tabularx}{\textwidth}{>{\raggedright\arraybackslash}p{4.4cm} >{\raggedright\arraybackslash}p{5.0cm} >{\raggedright\arraybackslash}X}
\headrow \thead{What a reviewer asks} & \thead{What the two inquiries show} & \thead{What they do not show}\\
\midrule
Does the company's work reach people outside it? & Two people read it and acted without being asked & How many read it and did not write \\
Could the company recruit on award? & Interest exists before any recruiting has begun & That either person would accept an offer \\
Is an eleven-role plan realistic? & Two inquiries arrived for a plan not yet funded & That either person fits a specific role \\
Would the team work inside the program's structure? & Both replies state the structured schedule first & Anything, until each person answers \\
\end{tabularx}
\end{apptable}
\vspace{-0.60cm}
\tabcap{Table 1. What two unprompted inquiries show a reviewer who scores the team,\\
set beside what they cannot show, so that no row can be read as a claim of hire.}

The right-hand column carries as much weight as the middle one. A company that
presents two inquiries as two hires has turned evidence into a claim, and a
reviewer who notices will discount every other number on the page.

Three features make the signal worth more than its size. It is unprompted, so
it carries no selection effect from the company's own outreach. It arrived
before an award, when most small businesses have only a recruiting timeline to
offer a reviewer. And it is the first test of the priorities the chief
executive set out on September 12, 2026, which the next subsection states as
terms of work.

\begin{appfloat}[!tb]
\begin{appfig}[mermaid-type / flowchart with two decisions and two stops]
% Two rows read as one path: left to right on the upper row, right to left on
% the lower row.  Each decision has a "no" branch that ends in a gray stop node,
% so the figure shows every place the thread can end, not only the happy path.
\node[mmin]   (a) at (0,0)     {LinkedIn message, no posting};
\node[mmstep] (b) at (3.65,0)  {Same reply to both: no open role today};
\node[mmdec,text width=19mm] (c) at (7.35,0) {Resume and consent to be named?};
\node[mmstep] (d) at (11.05,0) {One-page letter of intent, contingent on award};
\node[mmstep] (e) at (11.05,-2.35) {Named in an NIH or NSF application};
\node[mmdec,text width=19mm] (f) at (7.35,-2.35) {Award made?};
\node[mmstep] (g) at (3.65,-2.35) {Structured interview, six questions};
\node[mmgoal] (h) at (0,-2.35)  {Offer with pay range; I-9, EDD, DE 34};
\node[mmgray] (s1) at (7.35,1.75) {Thread closes courteously; nothing is kept};
\node[mmgray] (s2) at (7.35,-4.10) {No hire; the payroll earmark is released};
\draw[mmedge]  (a) -- (b);
\draw[mmedge]  (b) -- (c);
\draw[mmedgeb] (c) -- node[mmlabel,above] {yes} (d);
\draw[mmedged] (c) -- node[mmlabel,right] {no} (s1);
\draw[mmedge]  (d) -- (e);
\draw[mmedge]  (e) -- (f);
\draw[mmedgeb] (f) -- node[mmlabel,above] {yes} (g);
\draw[mmedged] (f) -- node[mmlabel,right] {no} (s2);
\draw[mmedgeb] (g) -- (h);
\ptitle{-1.55}{1.75}{From a message to a funded offer}
\pnote{-1.55}{-4.10}{Every node above the lower row\\is free of obligation to both sides}
\end{appfig}
\vspace{-0.60cm}
\figcaption{Figure 1. The path from an unprompted LinkedIn message to a funded offer, with the\\
two places the thread ends and the one point after which an obligation begins.}
\end{appfloat}

\subsection{The three priorities, stated as terms of work}

On September 12, 2026 the chief executive wrote to the President that the
startup's priorities had changed to three: obedience over hard work and
responsibilities; adaptation to existing authority, order, and structure; and a
teamwork-based, early-rise, scheduled routine. Day 3 of this block answers the
White House correspondence that followed. Day 1 applies the three priorities to
the only people who have so far asked to work inside them.

Each reply says, before anything else about the role, that the team works a
structured, early-start weekday schedule under established clinical and
regulatory procedures. The interview plan in \S4 tests the first two priorities
with two of its six questions: one asks for a time the candidate followed a
procedure they disagreed with, and the other asks who decides when a procedure
in a regulated study may be changed. A candidate who answers that the procedure's
owner and the reviewing body decide has understood the program's structure
before reading its protocol.
